\subsection{Survey of epidemic modeling software}\label{sec:methods-survey}

To place \texttt{epymorph} among existing software (Table~\ref{tab:comparison}), we conducted a structured, reproducible search for released packages that simulate or fit mechanistic epidemic models.
We describe it as a structured search rather than a systematic review, because the aim was an auditable and repeatable survey of software rather than a formal evidence synthesis.
The full protocol, the search string, and a summary of every package identified are given in Supplementary Section S1.

We drew packages from two kinds of sources.
First, we browsed curated software registries (the CRAN Task View on Epidemiology, rOpenSci, Bioconductor, PyPI, and the Journal of Open Source Software), searched Google Scholar and GitHub, and followed citations backward and forward from three anchor papers~\citep{king_statistical_2016,asfaw_tutorial_2024,lemaitre_flepimop_2024}.
Second, we searched two bibliographic databases, PubMed and Lens.org, for records from 2020 onward.
We used Lens.org because it indexes preprints, conference proceedings, and the Journal of Open Source Software, where software is often first announced.
The query required terms from three groups: the domain (e.g., epidemic, compartmental, metapopulation), the inference method (e.g., particle filter, ensemble Kalman, Bayesian, calibration), and the software itself (e.g., package, library, toolkit).

A package was included if it was released and installable, simulated or fit mechanistic compartmental or state-space epidemic models, provided at least one built-in method for Bayesian inference or uncertainty quantification, and had English documentation.
We excluded agent-based simulators without built-in inference, tools that only estimate the reproduction number, general-purpose statistical toolkits without epidemiological structure (e.g., Stan, PyMC), and packages with no activity in the preceding five years.

Together, the two databases yielded 3,033 unique records with abstracts.
Because this is too many to read in full, a keyword filter first advanced only the records whose title, abstract, or keywords mentioned software, inference, and epidemiology together, and this advanced 645 records to manual screening.
To measure what the filter missed, we audited the 2,388 records it set aside, reviewing every record whose title named a package (101 records) and a random sample of 100 of the rest, and we found no missed package that met our criteria.

The capabilities of each retained package were verified against its primary documentation (package website, CRAN or PyPI page, or software paper) as of September 2026, with a second author spot-checking the results.
In Table~\ref{tab:comparison}, a capability is marked present only if the package provides a dedicated construct, API, or data service for it; a capability that a user could program themselves does not count, because every package surveyed is programmable.
The survey retained 13 packages for comparison with \texttt{epymorph}: the 11 in Table~\ref{tab:comparison}, and two more (\texttt{LibBi} and \texttt{epidemia}) that are described in Supplementary Section S1.
The archived search exports, every screening decision, and the code that produces the counts above are provided with the manuscript's code.
